\section{The Arcade Learning Environment}
\label{sec:ale}

The previous section introduced two agents -- vanilla DQN and
Rainbow -- that differ in algorithmic sophistication but share the
same convolutional encoder and the same type of input: preprocessed
Atari game frames. This section introduces the benchmark on which
both agents are evaluated. The Arcade Learning
Environment~\citep{bellemare2013arcade} provides a standardized
interface to dozens of Atari 2600 games, making it possible to
evaluate a single agent architecture across a diverse set of tasks
without game-specific engineering. We describe the standard
evaluation protocol that established DQN's original results, the
data-efficient 100K variant used in our experiments, the metrics
used to compare agents across games, and the game-state access
that enables our interpretability analysis.

\subsection{The Standard Protocol}
\label{subsec:standard-protocol}

\citet{bellemare2013arcade} introduced the Arcade Learning
Environment (ALE): a software framework that wraps Atari 2600 game
ROMs into a standardized reinforcement learning interface. For
each game, the environment exposes $210 \times 160$ RGB frames
rendered at 60 frames per second, a discrete action set of up to
18 actions, a scalar reward signal derived from the game score, and
an episode termination signal. The key property is uniformity:
every game uses the same interface, so a single agent architecture
can be evaluated across the full suite without modification.

\smallskip
The games span a wide range of challenges -- spatial reasoning and
reaction time in Breakout, navigation and long-horizon planning in
Montezuma's Revenge, precise timing in Freeway, and multi-object
tracking in Space Invaders. This diversity is what made ALE the
standard benchmark for deep reinforcement learning: strong
performance across many games requires general-purpose
representations, not game-specific heuristics.

\smallskip
In the standard protocol, the agent does not act on every raw
frame. Instead, it selects an action every $k$-th frame, and the
chosen action is repeated for the intermediate frames.
\citet{mnih2015human} used $k = 4$, which became the default. This
\emph{frame skip} (or \emph{action repeat}) means that one agent
decision corresponds to four raw frames. The distinction between
raw frames and agent decisions matters when comparing training
budgets: 200 million raw frames is 50 million environment steps.
Frame skip serves two purposes: consecutive frames are nearly
identical, so acting on every one is unnecessary, and skipping
frames effectively speeds up the game from the agent's perspective.

\smallskip
The preprocessing pipeline that converts raw frames into the
$84 \times 84 \times 4$ input tensor was described in
\cref{subsec:dqn-architecture}. This pipeline -- max-pooling over
consecutive frames, grayscale conversion, resizing, and stacking
four frames -- is universal across methods. DQN, Rainbow, DrQ,
and SPR all receive the same preprocessed observations.

\smallskip
\citet{mnih2015human} trained DQN for 200 million frames per game,
roughly 38 days of game time at 60 frames per second. Each game
was trained independently with a single set of hyperparameters
shared across the entire suite. Evaluation followed a fixed
procedure: insert a random number (up to 30) of no-op actions at
the start of each episode to randomize the initial state, then let
the agent play for up to five minutes. Performance was averaged
over 30 evaluation episodes. Under this protocol, DQN achieved
human-competitive scores on roughly 60\% of 49 games -- the first
demonstration that a single architecture could learn diverse
visuomotor tasks from raw pixels.

\smallskip
\citet{machado2018revisiting} later identified weaknesses in the
original protocol and proposed two revisions that are now standard.
The first is \emph{sticky actions}: with probability 0.25, the
environment repeats the agent's previous action instead of
executing its current choice. Many Atari games are deterministic --
the same sequence of actions always produces the same outcome. An
agent trained on a deterministic game can exploit this by
memorizing fixed action sequences that achieve high scores without
learning general strategies. Sticky actions inject stochasticity
into the transition dynamics, forcing the agent to cope with
variability and preventing such memorization.

\smallskip
The second revision concerns episode boundaries. The original
protocol treated each life loss as an episode termination,
resetting the agent's trajectory. While this helped in some games
by making early deaths less costly to explore from, it changed the
effective task: the agent optimized for one-life performance rather
than full-game performance. \citet{machado2018revisiting}
recommended using only the game-over signal for episode boundaries,
preserving the original game structure. Both revisions -- sticky
actions and game-over termination -- are adopted in the 100K
protocol and in our experiments.

\smallskip
This protocol -- 200 million frames per game -- gave DQN enough
experience to learn effective policies for most games. But from a
practical standpoint, 200 million frames is an enormous amount of
interaction. The next section introduces a protocol that asks
whether agents can learn with far less.

\subsection{The 100K Protocol}
\label{subsec:100k-protocol}

\citet{kaiser2020model} proposed the Atari-100K benchmark: agents
are limited to 100{,}000 environment steps per game. With frame
skip $k = 4$, this corresponds to 400{,}000 raw frames -- roughly
six and a half minutes of game time at 60 frames per second.
Compared to the standard 200 million frames, the agent sees 0.2\%
of the data that DQN was originally trained on: a reduction of
500$\times$.

\smallskip
This constraint is not merely a smaller version of the standard
protocol. It creates a qualitatively different learning regime. In
the 200-million-frame setting, an agent can afford to explore
randomly for millions of steps, encounter diverse game situations
through sheer volume of experience, and gradually refine its
representations. At 100{,}000 steps, none of this is possible. The
agent must learn useful representations from its first experiences
and extract as much information as possible from every transition.
Vanilla DQN, which was designed for the high-data regime, barely
learns anything on most games under this constraint.

\smallskip
The benchmark has practical motivation beyond evaluating
algorithms. In real-world domains -- robotics, autonomous driving,
clinical decision support -- collecting millions of interactions is
expensive, slow, or dangerous. An agent that learns from 100{,}000
steps is closer to practical deployment than one requiring 200
million. The 100K benchmark distills this challenge into a
controlled setting where methods can be compared systematically.

\smallskip
The benchmark selects 26 games from the full 57-game ALE suite,
chosen to span a range of challenges: reward density, visual
complexity, planning horizon, and interaction patterns. Some games
offer frequent rewards that guide learning (Breakout, Pong), while
others require extended exploration before any reward is
encountered (Montezuma's Revenge, Private Eye). This diversity
ensures that strong aggregate performance requires general-purpose
data efficiency, not a strategy that works for only one type of
game.

\smallskip
After training for 100{,}000 steps, the agent is evaluated over
100 episodes using a near-greedy policy ($\varepsilon = 0.001$).
The evaluation score is the average undiscounted return across
these episodes. Sticky actions ($p = 0.25$) apply during both
training and evaluation, following the revised protocol of
\citet{machado2018revisiting}. The original benchmark paper
\citep{kaiser2020model} did not use sticky actions, but subsequent
work -- including SPR~\citep{schwarzer2021dataefficient} and the
Dopamine framework -- adopted them. Our experiments use sticky
actions consistently.

\smallskip
Hyperparameters tuned for 200 million frames are poorly suited to
the 100K regime. Key adjustments include a higher replay ratio --
performing more gradient updates per environment step to extract
more learning from each transition -- a shorter warmup period
before training begins, and faster $\varepsilon$ decay to shift
from exploration to exploitation earlier given the tiny budget.
\citet{kielak2020recent} demonstrated this concretely: the
model-free Rainbow baseline reported in the original benchmark
paper used default Dopamine hyperparameters without
game-specific tuning, making model-free methods appear weaker than
they are. A properly tuned Rainbow (``OTRainbow'') substantially
outperformed the original baseline without any algorithmic changes.
Published 100K results are sensitive to hyperparameter choices, not
just algorithmic differences -- a point that matters when comparing
across papers.

\smallskip
Comparing agents under the 100K constraint requires metrics that
aggregate performance across games with very different score
ranges. A score of 500 in one game might represent near-perfect
play, while the same score in another game is negligible. The next
section introduces the normalization and aggregation methods that
make these comparisons meaningful.

\subsection{Metrics}
\label{subsec:metrics}

Raw scores are not comparable across games. Pong scores range from
$-21$ to $+21$; Crazy Climber scores reach hundreds of thousands.
Averaging raw scores would let a single high-scoring game dominate
the aggregate, hiding poor performance elsewhere. Meaningful
comparison requires a normalization that puts all games on a common
scale.

\smallskip
The standard normalization is the \emph{human-normalized score}
(HNS):
\begin{equation}
\label{eq:hns}
\text{HNS} = \frac{\text{agent score} - \text{random score}}
                   {\text{human score} - \text{random score}}.
\end{equation}
An HNS of 0 means the agent performs at the level of a random
policy; an HNS of 1 means it matches human performance; values
above 1 are superhuman. In concrete terms, HNS measures what
fraction of the gap between random play and human play the agent
closes. The human and random baselines come from
\citet{mnih2015human}: human scores were recorded by a
professional games tester playing each game for two hours, and
random scores were measured from a uniform random policy.

\smallskip
Once scores are normalized, they must be aggregated across games
into a single number for comparison. The two obvious choices --
mean and median -- each have problems.
The mean is sensitive to outliers: if an agent achieves an HNS of
10 on one game (far above human level) but 0.1 on the rest, the
mean looks strong despite broadly weak performance. The median is
robust to outliers but statistically inefficient -- it uses only
the middle value and discards most of the data.
\citet{agarwal2021deep} demonstrated the practical consequences:
with only five runs per game (common in the literature), the median
of the same algorithm can vary by over 30\% depending on which
random seeds are included.

\smallskip
\citet{agarwal2021deep} recommended the \emph{interquartile mean}
(IQM) as the primary aggregate metric. IQM is a 25\% trimmed mean:
pool the normalized scores from all game-run combinations, discard
the bottom 25\% and top 25\%, and take the mean of the remaining
50\%. Discarding the top and bottom quartiles removes the outliers
that inflate the mean, while averaging the remaining half uses far
more of the data than the median's single central value. The result
is a metric that is more robust than the mean and more statistically
efficient than the median, producing tighter confidence intervals
for the same number of runs.

\smallskip
Point estimates of any aggregate metric are unreliable when the
number of runs is small. \citet{agarwal2021deep} proposed reporting
\emph{stratified bootstrap confidence intervals}: the resampling is
stratified by game -- for each game, the available runs are
resampled with replacement independently, so every game contributes
to every bootstrap sample rather than some games being over- or
under-represented by chance. The aggregate metric is computed on
each bootstrap sample, and the 95\% interval is reported. This
tells the reader how much the aggregate could shift with different
random seeds. A difference between two agents is meaningful only if
their confidence intervals show clear separation.

\smallskip
A complementary tool is the \emph{performance profile}: a curve
showing the fraction of game-run combinations where the agent
exceeds a given normalized score threshold. Higher curves indicate
stronger algorithms, and a curve that lies strictly above another
indicates stochastic dominance across the full range of thresholds
-- a stronger statement than a higher aggregate score, because it
means the better agent outperforms not just on average but at every
performance level. Performance profiles provide a richer picture
than any single
aggregate metric because they reveal the full distribution of
performance -- not just a summary statistic.

\smallskip
These metrics -- human-normalized scores aggregated via IQM with
bootstrap confidence intervals -- provide a principled framework
for comparing agents across games with very different score ranges.

\subsection{Game-State Access}
\label{subsec:game-state-access}

The metrics above measure performance, but to understand what the
agent's representations encode, we need to know what is actually
happening in the game. The agent sees the game as pixels, but the
emulator tracks the actual state: where each player is, what the score
is, how many lives remain.

\smallskip
The Atari 2600 stores this state in the console's random-access memory
(RAM), which holds just 128 values. Each value is a single number between 0 and 255, called a byte. Each
byte is stored at a fixed location in the memory, called an address. Each
address stores one game variable: one always holds the player's
horizontal position, another always holds the score. As the player moves
or the score changes, the number at that address changes with it. The agent never
sees these values, but we can read them directly from the emulator.

\smallskip
\citet{anand2019unsupervised} created the AtariARI benchmark by
manually identifying the game variables stored in RAM for 22
Atari games. For example, in Boxing, address 32 stores the player's x
position and address 34 stores the opponent's x position. In Pong,
address 49 stores the ball's y position. Each game has between 4 and 20 labeled variables, including
positions, scores, and lives.

\smallskip
To test whether the representation captures the ball's
position, we need to know where the ball actually is. RAM provides that
answer, and AtariARI's labels tell us which address to read. Without
these labels, we could measure how well the agent plays, but not what it
has learned to see. Linear probing
(\cref{subsec:representations}) uses these labeled values to test
exactly that.
